% ── docsmith handbook preamble — colours come from the generated _tokens.tex ──
% book class, Latin Modern serif body, navy palette, fancyhdr running heads,
% titlesec chapter/section styling, tcolorbox callout cards, lettrine, pullquote.

% PDF outline/bookmarks for navigation (screen-reader + reader wayfinding). Passed as
% options so they apply whenever pandoc's template loads hyperref, regardless of order.
\PassOptionsToPackage{bookmarks=true,bookmarksnumbered=true,bookmarksopen=true,pdfpagelabels=true}{hyperref}

\usepackage[paperwidth=6.5in,paperheight=9.25in,
            inner=0.9in,outer=0.75in,top=0.95in,bottom=1in,
            headsep=14pt,footskip=26pt]{geometry}

\usepackage{microtype}
\usepackage{graphicx}
\usepackage{enumitem}
\setlist{leftmargin=1.4em,topsep=3pt,itemsep=2pt,parsep=0pt}
\usepackage[most]{tcolorbox}
\usepackage{lettrine}
\usepackage{parskip}   % airy block paragraphs, matching the reference

\usepackage{titlesec}
% Chapter opener: large navy number flush right, navy rule, navy sans title.
\titleformat{\chapter}[display]
  {\sffamily\bfseries\color{navy}}
  {\raggedleft\fontsize{58}{58}\selectfont\thechapter}
  {6pt}
  {\color{navy}\titlerule[1.2pt]\vspace{10pt}\Huge\raggedright}
\titlespacing*{\chapter}{0pt}{6pt}{34pt}
% Front/back-matter titles (Contents, Terminology): serif bold black.
\titleformat{name=\chapter,numberless}[hang]
  {\rmfamily\bfseries\Huge\color{black}}{}{0pt}{}
\titlespacing*{name=\chapter,numberless}{0pt}{10pt}{24pt}

\titleformat{\section}
  {\sffamily\bfseries\large\color{navy}}{\thesection}{0.6em}{}
\titleformat{\subsection}
  {\sffamily\bfseries\normalsize\color{navysoft}}{\thesubsection}{0.5em}{}
\titlespacing*{\section}{0pt}{14pt}{6pt}
\titlespacing*{\subsection}{0pt}{10pt}{4pt}

\usepackage{fancyhdr}
\providecommand{\dslogo}{}
\providecommand{\dscompany}{}
\providecommand{\dsauthor}{}
% footer brand: logo (if supplied) + company name + author, muted grey.
% Renders inline as: [logo]  Company · Author  — the " · Author" half is
% suppressed when \dsauthor is empty, and the logo when \dslogo is empty.
\makeatletter
\newcommand{\dsfootbrand}{%
  {\def\dsnil{}\ifx\dslogo\dsnil\else
     \raisebox{-3.2pt}{\includegraphics[height=13pt]{\dslogo}}\hspace{7pt}\fi}%
  \small\sffamily\color{rulegrey}\dscompany
  {\edef\ds@a{\dsauthor}\ifx\ds@a\@empty\else\,\textperiodcentered\,\dsauthor\fi}}
\makeatother
\pagestyle{fancy}
\fancyhf{}
% L/R (not LE/LO/RE/RO) => identical placement on every page, one- or two-sided
\fancyhead[L]{\small\sffamily\color{rulegrey}\nouppercase{\leftmark}}
\fancyfoot[L]{\dsfootbrand}
\fancyfoot[R]{\small\sffamily\color{rulegrey}\thepage}
\renewcommand{\headrulewidth}{0.4pt}
\renewcommand{\footrulewidth}{0pt}
\renewcommand{\headrule}{\hbox to\headwidth{\color{rulegrey}\leaders\hrule height\headrulewidth\hfill}}
% chapter-opener pages use the 'plain' style — give them the SAME footer (logo +
% page number), no header rule, so numbering never jumps to the header/center.
\fancypagestyle{plain}{\fancyhf{}%
  \fancyfoot[L]{\dsfootbrand}\fancyfoot[R]{\small\sffamily\color{rulegrey}\thepage}%
  \renewcommand{\headrulewidth}{0pt}\renewcommand{\footrulewidth}{0pt}}

% ── Callout cards — tab label overlapping the top-left edge, rounded, shadow ──
\newtcolorbox{protip}{enhanced, breakable, boxrule=0.4pt, arc=2.4mm,
  colback=yellow!7!white, colframe=amber!70!white, drop fuzzy shadow=black!28,
  left=3mm,right=3mm,top=2.2mm,bottom=2.2mm,
  attach boxed title to top left={yshift=-2.6mm, xshift=4mm},
  boxed title style={colback=amber, colframe=amber, arc=1mm, boxrule=0pt},
  coltitle=black, title=Pro Tip, fonttitle=\bfseries\sffamily\footnotesize,
  before skip=10pt, after skip=8pt}

\newtcolorbox{anchorbox}{enhanced, breakable, boxrule=0.4pt, arc=2.4mm,
  colback=navytint, colframe=navysoft!55!white, drop fuzzy shadow=black!28,
  left=3mm,right=3mm,top=2.2mm,bottom=2.2mm,
  attach boxed title to top left={yshift=-2.6mm, xshift=4mm},
  boxed title style={colback=navy, colframe=navy, arc=1mm, boxrule=0pt},
  coltitle=white, title=Note, fonttitle=\bfseries\sffamily\footnotesize,
  before skip=10pt, after skip=8pt}

\newtcolorbox{warningbox}{enhanced, breakable, boxrule=0.4pt, arc=2.4mm,
  colback=orange!7!white, colframe=amberdeep!60!white, drop fuzzy shadow=black!28,
  left=3mm,right=3mm,top=2.2mm,bottom=2.2mm,
  attach boxed title to top left={yshift=-2.6mm, xshift=4mm},
  boxed title style={colback=amberdeep, colframe=amberdeep, arc=1mm, boxrule=0pt},
  coltitle=white, title=Watch Out, fonttitle=\bfseries\sffamily\footnotesize,
  before skip=10pt, after skip=8pt}

\newtcolorbox{plainbox}{enhanced, breakable, boxrule=0.4pt, arc=2.4mm,
  colback=violet!5!white, colframe=violet!45!white, drop fuzzy shadow=black!28,
  left=3mm,right=3mm,top=2.2mm,bottom=2.2mm,
  attach boxed title to top left={yshift=-2.6mm, xshift=4mm},
  boxed title style={colback=violet, colframe=violet, arc=1mm, boxrule=0pt},
  coltitle=white, title=In Plain English, fonttitle=\bfseries\sffamily\footnotesize,
  before skip=10pt, after skip=8pt}

\newtcolorbox{dobox}{enhanced, breakable, boxrule=0.4pt, arc=2.4mm,
  colback=green!5!white, colframe=greenok!50!white, drop fuzzy shadow=black!22,
  left=3mm,right=3mm,top=2.2mm,bottom=2.2mm,
  attach boxed title to top left={yshift=-2.6mm, xshift=4mm},
  boxed title style={colback=greenok, colframe=greenok, arc=1mm, boxrule=0pt},
  coltitle=white, title=Do, fonttitle=\bfseries\sffamily\footnotesize,
  before skip=10pt, after skip=4pt}

\newtcolorbox{dontbox}{enhanced, breakable, boxrule=0.4pt, arc=2.4mm,
  colback=red!4!white, colframe=amberdeep!45!white, drop fuzzy shadow=black!22,
  left=3mm,right=3mm,top=2.2mm,bottom=2.2mm,
  attach boxed title to top left={yshift=-2.6mm, xshift=4mm},
  boxed title style={colback=amberdeep, colframe=amberdeep, arc=1mm, boxrule=0pt},
  coltitle=white, title=Don't, fonttitle=\bfseries\sffamily\footnotesize,
  before skip=6pt, after skip=8pt}

% Cheatsheet — full-width red banner header, centred caps title, no shadow.
\newtcolorbox{cheatsheetbox}{enhanced, breakable, boxrule=0.6pt, arc=0mm,
  colback=white, colframe=cheatred, coltitle=white,
  colbacktitle=cheatred, halign title=center, toptitle=1mm, bottomtitle=1mm,
  left=3mm,right=3mm,top=2.2mm,bottom=2.2mm,
  title=\MakeUppercase{Cheatsheet}, fonttitle=\bfseries\sffamily\small,
  before skip=10pt, after skip=8pt}

% Pull quote — large navy italic with a thick navy left rule.
\newtcolorbox{pullquote}{enhanced, breakable, boxrule=0pt, sharp corners,
  colback=white, colframe=white, left=6mm, right=2mm, top=2mm, bottom=2mm,
  borderline west={2.5pt}{0pt}{navy},
  fontupper=\itshape\large\color{navy}, before skip=12pt, after skip=10pt}

% figure helper used by the diagram filter
\newcommand{\dgram}[3]{%
  \begin{center}\includegraphics[#2]{#1}\\[3pt]
  {\footnotesize\itshape\color{rulegrey}#3}\end{center}}

\usepackage{fvextra}
\fvset{fontsize=\footnotesize,breaklines=true,breakanywhere=true}
\providecommand{\tightlist}{\setlength{\itemsep}{1pt}\setlength{\parskip}{0pt}}

% ── Block code ──────────────────────────────────────────────────────────────
% Pandoc wraps every fenced code block (```bash, ```text, …) in a `Shaded`
% environment around an fvextra `Highlighting` verbatim. By default `Shaded` is a
% no-op, so code renders as bare text. Style it as a tinted card with a navy left
% rule so code blocks read like the callout cards. Guarded so it works whether or
% not pandoc has already defined `Shaded` (i.e. whether or not the doc has code).
\definecolor{codebg}{HTML}{F5F7FA}
\makeatletter
\@ifundefined{Shaded}{\newenvironment}{\renewenvironment}{Shaded}
  {\begin{tcolorbox}[breakable, enhanced, boxrule=0.6pt, arc=1.6mm,
     colback=codebg, colframe=navysoft!45!white,
     left=3.5mm, right=3mm, top=2mm, bottom=2mm,
     before skip=9pt, after skip=9pt]}
  {\end{tcolorbox}}
\makeatother
